\documentstyle[%


righttag%


,11pt%


,amssymb%


,russian%               remove in Eng.


%,izvru%                 Set izven% in Eng.


,izvru-ks%             for brief communications


]{amsart}





%


\newcommand{\E}{\bold E}


\newcommand{\N}{\bold N}


%% \def\RRp{{\rm I\kern-.25em R}_{+}}


%% \newcommand{RRp}{\Bbb R_+}


%% \def\IND{{\rm 1\kern-.25em I}}


\newcommand{\IND}{\Bbb I}





\newcommand{\tts}[1]{}





%\hoffset=-15mm%                  For Eng.


\setcounter{page}{74}


\god{2005}


\nomer{8~(519)} %                        Remove in Eng.


%\nomer{Vol.\,49, No.\,8}%               Set in Eng.


%\str{\pageref{begin}--\pageref{end}}%  Set in Eng.


\udk{519.216}





\author{\it З.A.\,Еникеева}


% NB:   ^^ remove in Eng.


\title{Принцип инвариантности для сумм \\ независимых 


случайных величин с замещениями\\ и модели финансового рынка}


\thanks{}





\date{}


%\pagestyle{myheadings}%         Set in Eng.


\pagestyle{plain} %              Remove in Eng.


\begin{document}


%\thispagestyle{plain}%          Set in Eng.


\maketitle


\label{begin}








В работе доказывается сходимость случайных ломаных


к винеровскому процессу с измененным временем


в пространстве непрерывных ограниченных функций


на действительной полупрямой.


Случайные ломаные определены суммами


независимых случайныx величин, причем в этих суммах


одни случайные величины случайным образом замещаются


другими. Замещения определяются умножением на значения


индикаторов, определенных на другом вероятностном пространстве,


элементы которого рассматриваются как случайные параметры,


и сходимость случайных ломаных доказывается для почти всех (п.\,в.)


значений этого случайного параметра.


Результаты данной работы являются обобщениями


соответствующих теорем из работ [1] и [2], полученных для сумм одинаково


распределенных случайных величин. Здесь рассматривается более


общая схема серий разнораспределенных случайных величин,


удовлетворяющих условию Линдеберга.





Заметим, что большое количество статей (см., напр., [3]--[5])


посвящено изучению сходимости


случайных ломаных, и данная работа принадлежит этому направлению.





Полученные предельные теоремы применяются в некоторых моделях стохастической


финансовой математики, в частности, приводится аналог формулы Блэка--Шольца


справедливой цены европейского опциона покупателя.








\section{Предельные теоремы}





Пусть $\{\Omega_1,\frak A_1 ,\bold P_1\}$ --- вероятностное пространство,


${\bold E}_1$ --- математическое ожидание относительно вероятности


${\bold P}_1$ и


пусть $A_{ji}(n)\in\frak A_1$ --- независимые события для любого $n\in\N$.


Будем предполагать, что вероятности этих событий не зависят от индекса $i$ и


$\bold P_1(A_{ji}(n))=a_{jn}$ для любого $i\in\N$.


Обозначим через $\IND_{ji}(n)(\omega_1)$ индикаторы событий


$A_{ji}(n)$, через $[b]$ и $\{b\}$ --- целую и дробную части числа $b$


соответственно.





Пусть $t\in \Bbb R_+ $. Определим последовательность функций следующим образом:


$f_n(t)=1$, если $[k_nt]=0$ и


$ f_n(t)=\E_1\,\IND_{[k_nt]i}(n) \IND_{([k_nt]-1)i}(n)\dotsc \IND_{1i}(n)


=a_{[k_nt]n}\dotsc a_{1n},$ если $[k_nt]>0$.





Пусть события $A_{ij}(n)$ удовлетворяют условию


\begin{itemize}


\item[(1)] существует функция $f(t)\in C_{b}(\Bbb R_+)$ такая, что


$f(t)= \lim \limits_{n\rightarrow \infty } f_{n}(t)$ для любого $t\in\Bbb R_+ $,


где $C_{b}(\Bbb R_+)$ --- пространство непрерывных ограниченных функций, 


заданных на $\Bbb R_+$.


\end{itemize}





Пусть $\{\Omega, \frak A, \bold P\}$ --- еще одно вероятностное пространство.


Через $\E$ будем обозначать математическое ожидание относительно 


вероятности $\bold P$.





Пусть $k_n\in \N$ такие, что $k_n<k_{n+1}$ для любого $n\in\N$ и $l\in\{1,2\}$.


Пусть $Y^{(l)}_{ni}$, $1\leq i \leq k_n$, --- последовательности серий 


независимых в каждой серии случайных величин, определенных на


$(\Omega, A, \bold P)$, такие, что для любого $n\in\N$ семейства


$\{Y^{(1)}_{ni}$, $1\leq i \leq k_n\}$ и


$\{Y^{(2)}_{ni}$, $1\leq i \leq k_n\}$ независимы.


Будем предполагать, что для $l\in\{1,2\}$ выполнены следующие условия:





\begin{itemize}


\item[(2)]  $\E\,Y^{(l)}_{ni} = 0$ для любых $n,i \in\N$;


\item[(3)] $\sum\limits_{i=1}^{k_n}\E\,(Y^{(l)}_{ni})^2 \rightarrow


v_l^2$ при $n\rightarrow\infty$ для некоторого $v_l^2\in \Bbb R_+$;


\item[(4)] $\sum\limits_{i=1}^{k_n}\E\,(Y^{(l)}_{ni})^2


\IND_{\{|Y^{(l)}_{ni}|>\tau\}}\rightarrow 0$


при $n\rightarrow\infty$ для любого $\tau\in\Bbb R_+$;


\item[(5)] $\sum\limits_{n=1}^{\infty}


\exp\big[-\frac{\varepsilon}{\max\limits_{1\le i\le k_n}\E


(Y^{(l)}_{ni})^2}\big]<\infty$ для любого $\varepsilon>0$.


\end{itemize}





Построим последовательность сумм


\begin{equation}


\begin{split}


S_0^{(n)}&=\sum_{i=1}^{k_n}Y_{ni}^{(1)},\\


%


S_1^{(n)}&=\sum_{i=1}^{k_n}\IND_{1i}(n)Y_{ni}^{(1)}+


\sum_{i=1}^{k_n}(1-\IND_{1i}(n))Y^{(2)}_{ni}, \\


%


&\dotsc\dotsc\dotsc \\


%


S_{k}^{(n)}&=\sum_{i=1}^{k_n}\IND_{1i}(n)\IND_{2i}(n)\dotsb


\IND_{ki}(n)Y_{ni}^{(1)}+


\sum_{i=1}^{k_n}(1-\IND_{1i}(n)\IND_{2i}(n)\dotsb


\IND_{ki}(n))Y^{(2)}_{ni}.


\end{split}


\tag{$\mathrm S$}


\end{equation}


\medskip





Рассмотрим последовательность случайных процессов


\begin{equation}


X_{n}(t)\equiv X_{n}(t;\omega_1)


\triangleq S^{(n)}_{[k_nt]}+\{k_nt\}( S^{(n)}_{[k_nt]+1} -


S^{(n)}_{[k_nt]}),\quad n\in\N.


\tag{$\mathrm X$}


\end{equation}





Пусть $W(t)$ и $W'(t)$ --- независимые винеровские процессы. Обозначим


$$


W_{fv_1v_2}(t)=W(v_1^2f(t))+W'(v_2^2(1-f(t))),\,\,\, t \in \Bbb R_+.


$$








\begin{thm}


Если выполнены условия {\em (1)--(5)} и


$X_{n}(\omega_1)$ --- последовательность


случайных процессов, определенная формулами $(\mathrm S)$ и $(\mathrm X)$, то


$X_{n}(\omega_1) @>d\;>> W_{fv_1v_2}$ при $n\to\infty$


в пространстве $C_b(\Bbb R_+)$ для п.\,в. $\omega_1 \in \Omega_1$.


\end{thm}





Обозначим


\begin{equation}


Z_{n}(x)\equiv Z_{n}(x;\omega_1)(Y)


=\frac {1}{k_n}\sum_{j=0}^{[k_nx]}S^{(n)}_j, \quad x\in[0,T],


\quad n\in\N,\quad \omega_1\in\Omega_1.


\tag{$\mathrm Z$}


\end{equation}








\begin{thm}


Если выполнены условия {\em (1)--(5)} и


$Z_{n}(\omega_1)$ --- последовательность


случайных процессов, определенная формулой $(\mathrm Z)$, то


$Z_n(\omega_1) @>d\;>> W''$ при $n\to\infty$


в пространстве $C[0,T]$ для п.\,в. $\omega_1 \in \Omega_1$, где


$W''$ --- центрированный гауссовский случайный процесс с функцией ковариации


$$


B(x_1, x_2)= v^2_1\bigg(2\int_0^{x_1} tf(t)dt+x_1\int_{x_1}^{x_2}


f(t)dt\bigg)+v^2_2\bigg(x_1x_2+2\int_0^{x_1} tf(t)dt-


(x_1+x_2)\int_0^{x_1}f(t)dt \bigg).


$$


\end{thm}








\section{Приложения к финансовой математике}





В работе [1] предложены три модели финансового рынка. 


В них действия агентов рассматриваются


как одинаково распределенные случайные величины. 


В данной работе будем предполагать, что


действия агентов --- разнораспределенные случайные величины.


Рассмотрим модель рынка с $n$ агентами, где $n$ достаточно велико.


Допустим, что торги ценными бумагами, которые в дальнейшем для 


простоты будем называть акциями, могут осуществляться только в


дискретные моменты времени $j\in\{0, 1,2,\dotsc, M \}$.


Предполагаем, что за период времени $\{0,1,2,\dotsc,M\}$ 


агент может только один раз


участвовать в торгах (купить или продать акции).


Если агент покупает акции, то их цена умножается на некоторое число, 


большее единицы. Если агент продает акции, то их цена умножается 


на число, меньшее единицы. Значит,


торговая политика $i$-го агента предполагается


биномиальной случайной величиной


$$


\xi_{i}^{(l)}=


\begin{cases}


U_{li}&


\text{с вероятностью } p_{li},\\


D_{li}& \text{с вероятностью } q_{li},\quad 1\le i\le n,\ \ l\in\{1,2\}.


\end{cases}


$$


Индекс $l$ введен для того, чтобы указать, что в зависимости от внешних 


обстоятельств политика агента может меняться, т.\,е. она может


быть случайной величиной $\xi_{i}^{(1)}$ или $\xi_{i}^{(2)}$.





Введем параметры $a_{li}$ и $v^2_{li}>0$, определив их равенствами


\begin{equation}\label{1}


\begin{split}


v^2_{li}&=(U_{li}-D_{li})^2p_{li}q_{li},\\


a_{li}v^2_{li} +1 &=U_{li}p_{li} + D_{li}q_{li},


\quad 1\le i\le n,\ \ l\in\{1,2\}.


\end{split}


\end{equation}





Пусть для $l\in\{1,2\}$ выполнены условия


\begin{itemize}


\item[$(1')$]$ \sum\limits_{i=1}^{n}v^2_{li} \rightarrow v^2_l$


при $n\rightarrow\infty$ для некоторого $v^2_l\in \Bbb R_+$;


\item[$(2')$] $\sum\limits_{i=1}^{n}\E(\ln \xi_{i}^{(l)}-\E\ln\xi_{i}^{(l)})^2


\IND_{\{|\ln \xi_{i}^{(l)}-\E\ln\xi_{i}^{(l)}|>\tau\}}\rightarrow 0$


при $n\rightarrow\infty$ для любого $\tau \in \Bbb R_+$;


\item[$(3')$] $\sum\limits_{n=1}^{\infty}


\exp\big[-\frac{\varepsilon}{\max\limits_{1 \le i\le n}v^2_{li}}\big]<\infty$


для любого $\varepsilon>0$ .


\end{itemize}





Параметр $a_{li}$ можно рассматривать, как относительный риск,


характеризующий разницу между ``рыночной'' мерой $(p_{li}, q_{li})$ и


``риск-нейтральной'' мерой $(\wt p_{li}, \wt q_{li})$, где числа


$\wt p_{li}$ и $\wt q_{li}$ определены равенствами


\begin{align*}  %%%%%%%% can be "gather"ed


&U\wt p_{li} + D\wt q_{li}=1,\\


&1 - \wt p_{li} =\wt q_{li}, \quad 1\le i\le n,\ \ l=1,2.


\end{align*}





Будем предполагать, что


\begin{itemize}


\item[$(4')$] $\sum\limits_{i=1}^{n} v^2_{li}a_{li}\to\mu_l-r_l$


при $n\to\infty$, $\mu_l-r_l\in\Bbb R_+$, $r_l>0$,


\end{itemize}


где $\mu_l$ --- ожидаемая средняя прибыль, $r_l$ --- банковский процент.





Рассмотрим случай, когда количество агентов на рынке остается неизменным, 


но агент может поменять свою тактику.


Тогда случайная величина $\xi_i^{(1)}$ заменяется на


$\xi_{i}^{(2)}.$


Итак, предполагаем, что выполнено следующее условие.


\begin{itemize}


\item[$(5')$] Действие любого агента не зависит от действий других агентов.


Случайная величина $\xi_{ik}$, описывающая действие $i$-го агента,


в $k$-й момент времени определяется следующим


образом:


\begin{itemize}


\item[(A)] если $\IND_{i1}^{(n)}(\omega_1)=1$,


$\IND_{i2}^{(n)}(\omega_1)=1,\dotsc,


\IND_{ij}^{(n)}(\omega_1)=1$,


то


$\xi_{ik}=\xi_i^{(1)}$, $k\le j$;


\item[(B)] если $\IND_{i1}^{(n)}(\omega_1)=1$,


$\IND_{i2}^{(n)}(\omega_1)=1,\dotsc,


\IND_{i(j-1)}^{(n)}(\omega_1)=1$, $\IND_{ij}^{(n)}(\omega_1)=0$,


то


$\xi_{ik}=\xi_{i}^{(2)}$, $k\ge j$.


\end{itemize}


Семейства случайных величин $\{\xi_i^{(1)}$, $1\leq i\leq n\}$


и $\{\xi_i^{(2)}$, $1\leq i\leq n\}$ независимы.


\end{itemize}





Построим случайный процесс, определяющий цену акций в момент 


времени $t\in [0, T]$


$$


S_n(t)=


S_0\bigg(\prod_{i=1}^{n}\prod_{j=0}^{[nt]}\xi_{ij}\bigg)^{1/n},


$$


где $S_0$ --- цена акции в начальный момент времени.








\begin{thm}


Если выполнены условия $(1)$ и $(1')$--$(5')$, то


$S_n(\omega_1)@>d\;>> S$ при $n\to\infty$ в пространстве


$L^{\infty}[0, T]$ для п.\,в. $\omega_1 \in \Omega_1$, где


$$


S(t)=S_0\exp\bigg\{{\sum_{l=1}^{2}(-1)^{l+1}


\bigg({\mu_l-r_l-\frac{v^2_l}{2}} \bigg)\int_{0}^{t} f(x)dx+


W''(t)+\bigg({\mu_2-r_2- \frac{v^2_2}{2}} \bigg)t} \bigg\},


$$


$t\in[0,T]$ и $W''$ --- гауссовский случайный процесс с функцией ковариации


$B(t_1,t_2)$, $t_1, t_2\in [0,T]$.


\end{thm}





Обозначим через $\wt S(t)$ случайный процесс, построенный точно так же,


как и $S_n(t)$, но


в котором случайные величины $\xi_{ij}$ обладают свойством


$\E\,\xi_{ij}=1$, $1\le i\le n$, $0\le j\le M$.








\begin{thm}


Если выполнены условия $(1)$ и $(1')$--$(5')$, то


$\wt S_n(\omega_1)@>d\;>> \wt S$ при $n\to\infty$


в пространстве $L^{\infty}[0, T]$ для п.\,в. $\omega_1 \in \Omega_1$, где


$$


\wt S(t)=S_0\exp\bigg\{ {\frac{v^2_2-v^2_1}{2}\int_{0}^{t} f(x)dx+


W''(t)- \frac{v^2_2}{2}t} \bigg\} ,\quad t\in[0,T].


$$


\end{thm}





Напомним, что справедливая цена $C_T$ европейского опциона покупателя


с датой исполнения $T$ и ценой покупки $K$ определяется формулой


Блэка--Шольца


$$


C_T={\E}\{ {e^{-rT}(\wt{S}(T) -K)_{+}}\} ,


$$


где случайный процесс $\wt S(t)$, $t\in [0, T]$,


является предельным по распределению для


последовательности $\wt S_n$ при $n\to\infty$.








\begin{thm}[{\series{m}\shape{n}\selectfont формула Блэка--Шольца}]


Предположим, что выполнены условия $(1)$ и\linebreak $(1')$--$(5')$.


Обозначим


$$


w^2(t)={\bold E}(W'(t))^2=B(t,t),\quad t\in [0,T].


$$


Тогда справедливая цена европейского опциона равна


$$


C_T=S_0\exp\bigg\{{\frac{w^2(T)}{2}-rT+\frac{v^2_2-v^2_1}{2}


\int_{0}^{T} f(x)dx-\frac{v_2^2T}{2}} \bigg\}


\Phi(\rho + w(T))-K\exp\{ {-rT}\} \Phi(\rho),


$$


где


$$


\rho=\bigg(\ln S_0-\ln K -\frac{v^2_1-v^2_1}{2}\int_0^T


f(x)dx-\frac{v_2^2T}{2}\bigg)\Big/ w(T).


$$


\end{thm}





При доказательстве теорем 3--5 существенно используется теорема 2.





\begin{thebibliography}{9}





\bibitem{1}


Chuprunov A.N., Rusakov O.V. {\it Convergence for


step line processes under summation of random indicators and


models of market pricing} // Lobachevskii J. of Math. -- Kazan, 2003.


\bibitem{2}


Fazekas I., Chuprunov A.N. {\it Convergence of random


step line to Ornstein--Uhlenbeck type processes} // Technical Report


of Debrecen University. -- 1996. -- \No\,24.


\bibitem{3}


Прохоров Ю.В. {\it Сходимость случайных процессов и предельные теоремы теории


вероятностей} // Теория вероятностей и ее применение. -- 1956.


-- T.\,1. -- \No\,2. -- C.\,176--238.


\bibitem{4}


Круглов В.М. {\it Слабая сходимость случайных ломаных к винеровскому процессу}


//  Теория вероятностей и ее применение. -- 1985. -- T.\,30. -- \No\,2.


-- C.\,208--218.


\bibitem{5}


Круглов В.М. {\it К принципу инвариантности}


// Теория вероятностей и ее применение. -- 1997. -- T.\,42.


-- \No\,2. -- C.\,239--261.








\end{thebibliography}





\vskip 2cm





\post{\it Казанский государственный}{\it Поступила}


\post{\it университет}{14.10.2004}





\label{end}


\end{document}








